\chapter{Tensor Algebra, Metrics, and Spacetime}

\section{The Minimum Is Technical}

The theoretical minimum is not a list of conclusions stated in words. It is
the least mathematics needed to move reliably to the next subject, which in
this case is general relativity and then cosmology. No advanced machinery is
required tonight, but the calculations must be followed patiently. Index
placement, transformation laws, and contractions are not decoration; they
are the grammar in which curved geometry is written.

Let \(x^{m}\) and \(y^{m}\) be two coordinate systems on the same space. A
covariant vector transforms as
\begin{equation}
  V_{m}^{(y)}(y)
  =
  \frac{\partial x^{n}}{\partial y^{m}}
  V_{n}^{(x)}(x),
  \label{eq:l04-covariant-law}
\end{equation}
whereas a contravariant vector transforms as
\begin{equation}
  W_{(y)}^{m}(y)
  =
  \frac{\partial y^{m}}{\partial x^{n}}
  W_{(x)}^{n}(x).
  \label{eq:l04-contravariant-law}
\end{equation}
The free index \(m\) appears on both sides and belongs to the \(y\)
description. The repeated \(n\) is summed and belongs to the \(x\)
description.

A collection of numbers is not a vector merely because it can be arranged
as an arrow. Temperature, pressure, and humidity are three scalars. Putting
their numerical values along three axes does not make a geometric vector,
because those values do not mix under a coordinate change according to
Eq.~\eqref{eq:l04-covariant-law} or
Eq.~\eqref{eq:l04-contravariant-law}. Transformation behavior defines the
object.

The two laws come directly from calculus. A scalar gradient satisfies
\begin{equation}
  \frac{\partial\phi}{\partial y^{m}}
  =
  \frac{\partial x^{n}}{\partial y^{m}}
  \frac{\partial\phi}{\partial x^{n}},
  \label{eq:l04-gradient-chain-rule}
\end{equation}
so gradients are covariant. Coordinate differentials satisfy
\begin{equation}
  dy^{m}
  =
  \frac{\partial y^{m}}{\partial x^{n}}\,dx^{n},
  \label{eq:l04-differential-chain-rule}
\end{equation}
so infinitesimal displacements are contravariant. These chain rules are the
archetypes from which the terminology is abstracted.

\begin{figure}[tbp]
  \centering
  \includegraphics[width=0.84\textwidth]{lecture_04_figure_01.png}
  \caption{The two chain-rule archetypes: a scalar gradient transforms
  covariantly, while a differential displacement transforms
  contravariantly.}
  \lectureframe{00:12:35}
\end{figure}

\begin{classroomqa}
\audiencequestion{If an object satisfies one vector transformation law,
must it also satisfy the other?
\lecturetimestamp{00:07:24--00:08:33}}
\lecturerresponse{No. Before a metric is introduced, a covector such as
\(\partial_{m}\phi\) and a vector such as \(dx^{m}\) are different kinds of
object. The same letter should not be used to identify them. The metric will
later provide a one-to-one map between their component descriptions.}
\end{classroomqa}

Coordinates conventionally carry upper indices. A coordinate in the
denominator of a derivative therefore contributes a lower index, while a
coordinate differential naturally carries an upper one. The convention is
arbitrary, but once adopted it makes the Jacobian patterns difficult to
misread.

\section{Higher-Rank Tensors and Contraction}

Each tensor index transforms independently. For a tensor with two lower
indices,
\begin{equation}
  T_{mr}^{(y)}
  =
  \frac{\partial x^{n}}{\partial y^{m}}
  \frac{\partial x^{s}}{\partial y^{r}}
  T_{ns}^{(x)}.
  \label{eq:l04-covariant-rank-two}
\end{equation}
For one upper and one lower index,
\begin{equation}
  \bigl(T^{m}{}_{r}\bigr)_{(y)}
  =
  \frac{\partial y^{m}}{\partial x^{n}}
  \frac{\partial x^{s}}{\partial y^{r}}
  \bigl(T^{n}{}_{s}\bigr)_{(x)}.
  \label{eq:l04-mixed-rank-two}
\end{equation}
The general rule is one direct Jacobian for every upper index and one
inverse Jacobian for every lower index.

\begin{figure}[tbp]
  \centering
  \includegraphics[width=0.86\textwidth]{lecture_04_figure_02.png}
  \caption{Covariant and mixed rank-two transformation laws. The completed
  board displays one Jacobian factor for each tensor index.}
  \lectureframe{00:17:45}
\end{figure}

\begin{classroomqa}
\audiencequestion{May upper and lower indices be mixed on the same tensor?
\lecturetimestamp{00:15:15--00:16:18}}
\lecturerresponse{Yes. A mixed tensor is perfectly legitimate. Each index
retains its own transformation rule, as
Eq.~\eqref{eq:l04-mixed-rank-two} shows.}
\end{classroomqa}

The next operation is \emph{contraction}. Identify one upper index with one
lower index and sum:
\begin{equation}
  T^{m}{}_{m}
  =
  T^{1}{}_{1}+T^{2}{}_{2}+\cdots.
  \label{eq:l04-trace}
\end{equation}
No free index remains, so the result should be a scalar. The transformation
law verifies it:
\begin{align}
  \bigl(T^{m}{}_{m}\bigr)_{(y)}
  &=
  \frac{\partial y^{m}}{\partial x^{r}}
  \frac{\partial x^{s}}{\partial y^{m}}
  \bigl(T^{r}{}_{s}\bigr)_{(x)}
  \nonumber\\
  &=
  \delta^{s}{}_{r}\,T^{r}{}_{s}{}^{(x)}
  =
  \bigl(T^{r}{}_{r}\bigr)_{(x)}.
  \label{eq:l04-trace-invariant}
\end{align}
The paired Jacobians are inverse matrices. Their product is the identity,
represented by the Kronecker delta.

\begin{figure}[tbp]
  \centering
  \includegraphics[width=0.86\textwidth]{lecture_04_figure_03.png}
  \caption{Contraction of a mixed tensor. The circled trace has no free
  indices and is unchanged by the coordinate transformation.}
  \lectureframe{00:22:35}
\end{figure}

The same argument works at any rank. Contracting one upper with one lower
index removes those two indices and leaves a tensor whose remaining indices
transform normally. A familiar special case is the inner product
\begin{equation}
  V^{m}W_{m},
  \label{eq:l04-inner-product}
\end{equation}
which is a scalar.

\begin{classroomqa}
\audiencequestion{Does contraction imply that the original tensor was a
scalar times the lower-rank tensor that remains?
\lecturetimestamp{00:25:00--00:26:50}}
\lecturerresponse{No. Contraction constructs a new object and generally
discards information. A special tensor may factor, but no such factorization
follows from contraction alone.}
\end{classroomqa}

\begin{classroomqa}
\audiencequestion{Does a two-index tensor necessarily represent a plane?
\lecturetimestamp{00:26:50--00:27:35}}
\lecturerresponse{Sometimes a rank-two object admits that geometric
interpretation, but rank by itself does not force it. The metric tensor is
the important example here: its two indices encode the bilinear rule for
distance.}
\end{classroomqa}

\section{The Metric from the Scalar Line Element}

For two neighboring points separated by \(dx^{m}\), define the squared
distance
\begin{equation}
  \boxed{
  ds^{2}
  =
  g_{mn}(x)\,dx^{m}dx^{n}.}
  \label{eq:l04-line-element}
\end{equation}
Both vector indices are contracted against the two lower metric indices, so
\(ds^{2}\) is a scalar.

\begin{figure}[tbp]
  \centering
  \includegraphics[width=0.82\textwidth]{lecture_04_figure_04.png}
  \caption{The metric tensor enters as the bilinear coefficient that turns
  an infinitesimal displacement into its squared length.}
  \lectureframe{00:28:35}
\end{figure}

To prove that \(g_{mn}\) is a tensor, rewrite the same scalar in the
\(y\)-coordinates:
\begin{equation}
  dx^{m}
  =
  \frac{\partial x^{m}}{\partial y^{r}}\,dy^{r}.
\end{equation}
Substitution gives
\begin{align}
  ds^{2}
  &=
  g_{mn}^{(x)}
  \frac{\partial x^{m}}{\partial y^{r}}
  \frac{\partial x^{n}}{\partial y^{s}}
  dy^{r}dy^{s}
  \nonumber\\
  &=
  g_{rs}^{(y)}\,dy^{r}dy^{s},
  \label{eq:l04-metric-rewrite}
\end{align}
and therefore
\begin{equation}
  \boxed{
  g_{rs}^{(y)}
  =
  \frac{\partial x^{m}}{\partial y^{r}}
  \frac{\partial x^{n}}{\partial y^{s}}
  g_{mn}^{(x)}.}
  \label{eq:l04-metric-transform}
\end{equation}
This is exactly the covariant rank-two law.

\begin{figure}[tbp]
  \centering
  \includegraphics[width=0.88\textwidth]{lecture_04_figure_05.png}
  \caption{The line element rewritten in \(y\)-coordinates. The boxed
  coefficient of \(dy^{r}dy^{s}\) is the transformed metric
  \(g_{rs}^{(y)}\).}
  \lectureframe{00:31:25}
\end{figure}

\begin{classroomqa}
\audiencequestion{Is \(ds^{2}\) really the same in every coordinate system?
\lecturetimestamp{00:45:25--00:46:25}}
\lecturerresponse{Yes. It is one scalar interval written in two coordinate
languages. Its invariance is not a result discovered after the
transformation law; it is the premise that forces
Eq.~\eqref{eq:l04-metric-transform}.}
\end{classroomqa}

\section{The Inverse Metric}

At a fixed point, the components \(g_{mn}\) form a matrix. The metric used
for geometry is symmetric and nondegenerate:
\begin{equation}
  g_{mn}=g_{nm},
  \qquad
  \det(g_{mn})\ne0.
  \label{eq:l04-metric-properties}
\end{equation}
It therefore has an inverse matrix. Denote the inverse components by
\(g^{mn}\). Matrix multiplication becomes
\begin{equation}
  \boxed{
  g^{mr}g_{rn}
  =
  \delta^{m}{}_{n}.}
  \label{eq:l04-inverse-metric}
\end{equation}
The inverse components transform as a contravariant rank-two tensor. Thus
\(g_{mn}\) and \(g^{mn}\) contain the same geometric information.

\begin{figure}[tbp]
  \centering
  \includegraphics[width=0.82\textwidth]{lecture_04_figure_06.png}
  \caption{Matrix inversion translated into tensor notation. The identity
  tensor has one upper and one lower index.}
  \lectureframe{00:36:55}
\end{figure}

\begin{classroomqa}
\audiencequestion{Should both indices on the Kronecker delta be lower?
\lecturetimestamp{00:33:20--00:34:32}}
\lecturerresponse{Not in the invariant identity-map notation used here.
\(\delta^{m}{}_{n}\) carries one upper and one lower index because it maps a
vector space to itself. In an orthonormal Cartesian calculation one often
writes \(\delta_{mn}\), but that restricted convenience should not be
confused with the general identity tensor.}
\end{classroomqa}

In polar coordinates on the plane,
\begin{equation}
  ds^{2}=dr^{2}+r^{2}d\theta^{2},
\end{equation}
so
\begin{equation}
  (g_{mn})
  =
  \begin{pmatrix}
    1&0\\
    0&r^{2}
  \end{pmatrix},
  \qquad
  (g^{mn})
  =
  \begin{pmatrix}
    1&0\\
    0&r^{-2}
  \end{pmatrix}.
  \label{eq:l04-polar-inverse}
\end{equation}
For a diagonal matrix, inversion simply reciprocates each diagonal entry.

\begin{classroomqa}
\audiencequestion{Is the angular metric component \(r\) or \(r^{2}\)?
\lecturetimestamp{00:37:30--00:38:24}}
\lecturerresponse{It is \(r^{2}\). The coordinate differential is already
\(d\theta\); the squared physical arc length is
\((r\,d\theta)^{2}=r^{2}d\theta^{2}\).}
\end{classroomqa}

\begin{classroomqa}
\audiencequestion{Can a tensor fail to have an inverse?
\lecturetimestamp{00:44:00--00:45:25}}
\lecturerresponse{A general rank-two tensor can be singular. A metric, by
definition, is nondegenerate and therefore has no zero eigenvalue. Symmetry
alone is not enough; nondegeneracy is the condition that guarantees the
inverse. The same statement holds for a Lorentzian metric despite its minus
signs.}
\end{classroomqa}

Only the symmetric part contributes to the quadratic form. Indeed,
\begin{equation}
  g_{12}\,dx^{1}dx^{2}
  +
  g_{21}\,dx^{2}dx^{1}
  =
  (g_{12}+g_{21})\,dx^{1}dx^{2}.
  \label{eq:l04-symmetric-part}
\end{equation}
An antisymmetric part would cancel because the coordinate differentials
commute.

\begin{figure}[tbp]
  \centering
  \includegraphics[width=0.82\textwidth]{lecture_04_figure_08.png}
  \caption{The two off-diagonal contributions combine through
  \(g_{12}+g_{21}\), showing why the line element depends only on the
  symmetric part of the metric.}
  \lectureframe{00:47:45}
\end{figure}

\begin{classroomqa}
\audiencequestion{Must \(g_{12}\) and \(g_{21}\) be equal?
\lecturetimestamp{00:46:30--00:49:00}}
\lecturerresponse{For the metric, we take them equal without loss of
geometric information. Any antisymmetric difference contributes nothing to
\(ds^{2}\), so retaining it would add redundant components rather than new
distance data.}
\end{classroomqa}

\section{Raising and Lowering Indices}

The metric now links the two vector descriptions:
\begin{equation}
  V_{n}=g_{nm}V^{m},
  \qquad
  V^{m}=g^{mn}V_{n}.
  \label{eq:l04-raise-lower-vector}
\end{equation}
The first equation lowers an index and the second raises it. Because the
metric is invertible, either set of components determines the other. In
geometric terms, \(V_{n}\) is the covector obtained by taking the metric
inner product of \(V\) with each coordinate basis vector. In an orthonormal
basis this reduces to the familiar projection picture.

\begin{figure}[tbp]
  \centering
  \includegraphics[width=0.82\textwidth]{lecture_04_figure_07.png}
  \caption{Lowering a vector index with \(g_{mn}\). The metric turns the
  contravariant components \(V^{m}\) into covariant components \(V_{n}\).}
  \lectureframe{00:43:15}
\end{figure}

Indices on higher-rank tensors are handled one at a time. For example,
\begin{equation}
  T^{m}{}_{n}
  =
  g^{mr}T_{rn},
  \label{eq:l04-raise-tensor-index}
\end{equation}
raises the first index of a covariant rank-two tensor.

\begin{figure}[tbp]
  \centering
  \includegraphics[width=0.84\textwidth]{lecture_04_figure_09.png}
  \caption{The vector inverse relation and the corresponding operation of
  raising one index on a rank-two tensor.}
  \lectureframe{00:52:30}
\end{figure}

Apply the same operation to the displacement:
\begin{equation}
  dx_{m}=g_{mn}\,dx^{n}.
  \label{eq:l04-lower-displacement}
\end{equation}
The squared interval then has three equivalent forms,
\begin{equation}
  \boxed{
  ds^{2}
  =
  g_{mn}\,dx^{m}dx^{n}
  =
  dx^{m}dx_{m}
  =
  g^{mn}\,dx_{m}dx_{n}.}
  \label{eq:l04-three-distance-forms}
\end{equation}

\begin{figure}[tbp]
  \centering
  \includegraphics[width=0.86\textwidth]{lecture_04_figure_10.png}
  \caption{The completed board comparison of the three contracted forms of
  the same scalar interval.}
  \lectureframe{00:58:15}
\end{figure}

\begin{classroomqa}
\audiencequestion{Was the first expression on the board meant as an equality
between \(dx^{m}dx_{n}\) and its contraction?
\lecturetimestamp{00:56:30--00:57:45}}
\lecturerresponse{No. That equals sign is withdrawn in the discussion. One
first identifies an upper and lower index and sums; only after carrying out
that contraction does one obtain the scalar
\(dx^{m}dx_{m}=ds^{2}\).}
\end{classroomqa}

\begin{classroomqa}
\audiencequestion{What is the relation between the upper- and lower-index
metrics in these formulas?
\lecturetimestamp{00:58:20--00:59:55}}
\lecturerresponse{They are inverse matrices, related by
Eq.~\eqref{eq:l04-inverse-metric}. For a diagonal metric the entries are
reciprocals. With off-diagonal components one performs an ordinary matrix
inverse.}
\end{classroomqa}

\section{Tensor Equations and the Deferred Calculus}

Tensors can be added only when their index types agree. The sum
\(T_{mn}+W_{mn}\) transforms coherently; a formal sum of \(T_{mn}\) and
\(W^{mn}\) does not.

If every component of a tensor vanishes in one coordinate system, every
component vanishes in all coordinate systems, because each transformed
component is a linear combination of zeros. Hence if
\begin{equation}
  T_{mn}=W_{mn}
\end{equation}
in one chart, then
\begin{equation}
  T_{mn}-W_{mn}=0
\end{equation}
there and therefore everywhere. Tensor equality is geometric even though
the separate components depend on coordinates.

\begin{figure}[tbp]
  \centering
  \includegraphics[width=0.84\textwidth]{lecture_04_figure_11.png}
  \caption{A tensor equality recast as the vanishing of the difference.
  Once the difference is the zero tensor in one chart, it is zero in every
  chart.}
  \lectureframe{01:02:40}
\end{figure}

This completes tensor algebra, not tensor calculus. Ordinary partial
derivatives of tensor components do not generally transform as tensor
components, because differentiating the Jacobians produces extra terms.
The corrected derivative will be introduced later. For now the lecture
deliberately turns from space to spacetime.

\section{From Cartesian Space to Proper Time}

Even flat Euclidean space has many Cartesian frames. They may be translated
or rotated relative to one another. In two dimensions,
\begin{equation}
  ds^{2}
  =
  (dx^{1})^{2}+(dx^{2})^{2}
  =
  \delta_{mn}\,dx^{m}dx^{n}.
  \label{eq:l04-euclidean-interval}
\end{equation}
Transformations between Cartesian frames preserve this Pythagorean form.
In that restricted setting the identity metric makes upper and lower
components numerically equal.

\begin{figure}[tbp]
  \centering
  \includegraphics[width=0.84\textwidth]{lecture_04_figure_12.png}
  \caption{Two Cartesian frames on the same flat plane together with the
  Euclidean interval they preserve.}
  \lectureframe{01:09:25}
\end{figure}

Spacetime has one time coordinate and three spatial coordinates. In one
space dimension, a Euclidean-looking sum \(dt^{2}+dx^{2}\) is not preserved
by Lorentz transformations. The invariant quantity is the proper-time
interval
\begin{equation}
  \boxed{
  d\tau^{2}
  =
  dt^{2}-dx^{2}}
  \qquad (c=1).
  \label{eq:l04-proper-time-1plus1}
\end{equation}

\begin{figure}[tbp]
  \centering
  \includegraphics[width=0.80\textwidth]{lecture_04_figure_13.png}
  \caption{Two nearby events on a worldline and the invariant proper-time
  interval \(d\tau^{2}=dt^{2}-dx^{2}\).}
  \lectureframe{01:13:35}
\end{figure}

Setting \(c=1\) makes time and distance use the same unit. Keeping different
units would be like measuring one Cartesian direction in feet and another
in inches: conversion factors would obscure the geometric symmetry. If
\(c\) is restored,
\begin{equation}
  d\tau^{2}
  =
  dt^{2}-\frac{dx^{2}}{c^{2}},
\end{equation}
when \(d\tau\) is measured in time units.

\begin{classroomqa}
\audiencequestion{How can particle physicists express length, time, and mass
in electron volts?
\lecturetimestamp{01:16:00--01:17:32}}
\lecturerresponse{Setting \(c=1\) identifies time units with length units.
Setting \(\hbar=1\) additionally relates inverse length and inverse time to
energy or mass. Only \(c=1\) is needed in the present relativistic
discussion; the \(\hbar=1\) convention belongs to quantum theory.}
\end{classroomqa}

In three spatial dimensions,
\begin{equation}
  \boxed{
  d\tau^{2}
  =
  dt^{2}-dx^{2}-dy^{2}-dz^{2}.}
  \label{eq:l04-proper-time-3plus1}
\end{equation}
Special relativity may be characterized as the theory of inertial
coordinate transformations that preserve this indefinite interval.

\begin{figure}[tbp]
  \centering
  \includegraphics[width=0.82\textwidth]{lecture_04_figure_14.png}
  \caption{The proper-time interval extended to one time and three spatial
  dimensions.}
  \lectureframe{01:18:40}
\end{figure}

\section{General Spacetime Coordinates}

Accelerated coordinates have already shown that curvilinear coordinates are
useful in spacetime even before genuine gravitational curvature is
introduced. Write
\begin{equation}
  x^{\mu}=(x^{0},x^{1},x^{2},x^{3}),
  \qquad
  x^{0}=t.
\end{equation}
Greek indices run over spacetime coordinates. In arbitrary coordinates the
proper-time element takes the tensorial form
\begin{equation}
  \boxed{
  d\tau^{2}
  =
  g_{\mu\nu}(x)\,dx^{\mu}dx^{\nu}.}
  \label{eq:l04-general-spacetime-metric}
\end{equation}

\begin{figure}[tbp]
  \centering
  \includegraphics[width=0.86\textwidth]{lecture_04_figure_15.png}
  \caption{The general spacetime line element and the \(x^{\mu}\) and
  \(y^{\mu}\) coordinate labels.}
  \lectureframe{01:22:45}
\end{figure}

\begin{classroomqa}
\audiencequestion{Why number time by \(0\), and may a spacetime have more
than one time coordinate?
\lecturetimestamp{01:20:00--01:21:50}}
\lecturerresponse{Starting with \(x^{0}=t\) is a convention that leaves
\(x^{1},x^{2},x^{3}\) for the spatial coordinates. Higher-dimensional
models may add spatial coordinates, but the Lorentzian geometry used here
has one timelike direction. Multiple-time theories are outside the present
framework.}
\end{classroomqa}

In Minkowski coordinates,
\begin{equation}
  \eta_{\mu\nu}
  =
  \operatorname{diag}(1,-1,-1,-1),
  \label{eq:l04-eta}
\end{equation}
and
\begin{equation}
  d\tau^{2}
  =
  \eta_{\mu\nu}\,dx^{\mu}dx^{\nu}.
  \label{eq:l04-minkowski-line-element}
\end{equation}
The symbol \(\eta_{\mu\nu}\) is the spacetime counterpart of the Cartesian
\(\delta_{mn}\), but its signs encode a Lorentzian rather than Euclidean
geometry.

\begin{figure}[tbp]
  \centering
  \includegraphics[width=0.84\textwidth]{lecture_04_figure_16.png}
  \caption{The Minkowski metric matrix with one positive and three negative
  diagonal entries.}
  \lectureframe{01:26:10}
\end{figure}

Under every nonsingular real coordinate transformation, the numbers of
positive and negative metric eigenvalues are preserved. With the convention
used here, a physical four-dimensional spacetime metric has one positive
and three negative directions.

\begin{classroomqa}
\audiencequestion{Could a coordinate change simply interchange the labels
\(x\) and \(t\)?
\lecturetimestamp{01:27:50--01:29:18}}
\lecturerresponse{Labels may be exchanged, and arbitrary smooth coordinates
may mix space and time. The metric components then change accordingly.
What cannot be changed by a regular coordinate transformation is the
Lorentzian signature: there remains one timelike and three spacelike
directions.}
\end{classroomqa}

\begin{classroomqa}
\audiencequestion{Is the Minkowski metric its own inverse?
\lecturetimestamp{01:29:18--01:30:00}}
\lecturerresponse{Yes. Each diagonal entry is \(+1\) or \(-1\), and its
reciprocal is itself, so
\(\eta^{\mu\nu}=\operatorname{diag}(1,-1,-1,-1)\).}
\end{classroomqa}

\section{A First Expanding-Spacetime Metric}

For a homogeneous and isotropic flat spatial geometry, choose comoving
coordinates and write
\begin{equation}
  (g_{\mu\nu})
  =
  \begin{pmatrix}
    1&0&0&0\\
    0&-a^{2}(t)&0&0\\
    0&0&-a^{2}(t)&0\\
    0&0&0&-a^{2}(t)
  \end{pmatrix}.
  \label{eq:l04-flat-flrw-matrix}
\end{equation}
The corresponding interval is
\begin{equation}
  \boxed{
  d\tau^{2}
  =
  dt^{2}
  -
  a^{2}(t)\left(dx^{2}+dy^{2}+dz^{2}\right).}
  \label{eq:l04-flat-flrw}
\end{equation}
The single function \(a(t)\) is the scale factor. At a fixed cosmic time,
two comoving points separated by coordinate distance \(\Delta\ell_{\rm com}\)
have physical distance
\begin{equation}
  \Delta\ell_{\rm phys}(t)
  =
  a(t)\,\Delta\ell_{\rm com}.
  \label{eq:l04-physical-distance}
\end{equation}
If \(a(t)\) grows, the physical separation grows even while the comoving
coordinates remain fixed.

\begin{figure}[tbp]
  \centering
  \includegraphics[width=0.86\textwidth]{lecture_04_figure_17.png}
  \caption{A homogeneous time-dependent metric with one common spatial
  scale factor controlling all three spatial directions.}
  \lectureframe{01:32:30}
\end{figure}

\begin{classroomqa}
\audiencequestion{What does it mean geometrically to say that the universe
is expanding?
\lecturetimestamp{01:30:00--01:33:00}}
\lecturerresponse{It means that the spatial part of the spacetime metric
changes with time. In comoving coordinates, fixed coordinate separations
correspond to physical distances proportional to \(a(t)\). Expansion is
therefore a statement about the metric, not motion through a preexisting
external space.}
\end{classroomqa}

Two-dimensional surfaces make the role of a changing metric coefficient
visible. A rotationally symmetric surface can be written
\begin{equation}
  ds^{2}
  =
  dr^{2}+f^{2}(r)\,d\theta^{2}.
  \label{eq:l04-surface-profile}
\end{equation}
The circumference of the circle at fixed \(r\) is \(2\pi f(r)\).

For the flat plane, \(f(r)=r\). For a unit sphere,
\begin{equation}
  f(r)=\sin r,
  \qquad
  ds^{2}=dr^{2}+\sin^{2}r\,d\theta^{2}.
  \label{eq:l04-sphere-analogy}
\end{equation}
The circles grow from one pole to the equator and then shrink toward the
other pole.

\begin{figure}[tbp]
  \centering
  \includegraphics[width=0.84\textwidth]{lecture_04_figure_18.png}
  \caption{The sphere and its metric. The factor \(\sin r\) controls the
  circumference of circles at fixed polar distance \(r\).}
  \lectureframe{01:35:30}
\end{figure}

For \(f(r)=e^{r}\),
\begin{equation}
  ds^{2}=dr^{2}+e^{2r}d\theta^{2},
  \label{eq:l04-horn-analogy}
\end{equation}
the circumference grows exponentially and the surface has a horn-like
profile. More complicated growing and oscillating choices of \(f(r)\)
produce correspondingly complicated surfaces.

\begin{figure}[tbp]
  \centering
  \includegraphics[width=0.84\textwidth]{lecture_04_figure_19.png}
  \caption{A general radial profile and a horn-like surface. The coefficient
  of \(d\theta^{2}\) determines how circumference changes along \(r\).}
  \lectureframe{01:37:35}
\end{figure}

These surfaces vary with position \(r\); the cosmological metric varies
with time \(t\). The analogy is not yet a curvature calculation. Its purpose
is to make metric coefficients readable as geometry. The next steps are
parallel transport, curvature, and eventually Einstein's field equations.
